\chapter{Conclusions}

This work developed an explainable reviewed-item retrieval system for restaurant search. The
main proposed configuration uses multilingual dense sentence embeddings and top-10 mean
aggregation to rank restaurants and return original review sentences as evidence.

On RIRD, the proposed model achieved nDCG@10 0.682 and MAP 0.618, outperforming the fair BM25
sentence baseline and approaching the supervised fine-tuned reference. The experiments also
showed that top-$k$ aggregation is more reliable than max or averaging every review, and that
review-volume correlations can be reduced by capping the number of reviews.

The B3 analysis produced an important negative result: applying query rewriting and a generic
cross-encoder to every query reduced the overall ranking score. The deployed pipeline therefore
uses the reranker selectively for explicit negation and exclusion queries, while preserving the
strong dense baseline for ordinary queries. This is a practical example of using a more
complex model only where its measured benefit justifies its cost.

The Thai component demonstrates sentence-level retrieval over real Wongnai reviews, but cannot
claim Thai restaurant ranking because restaurant identifiers are unavailable. Future work
should complete the human support and Thai annotation studies, evaluate larger restaurant-level
Thai datasets, and train or adapt a reranker on reviewed-item relevance. Overall, the project
shows that semantic retrieval becomes substantially more useful when paired with transparent
evidence, fair evaluation and explicit acknowledgement of its limitations.
